\chapter*{LỜI NÓI ĐẦU}
\addcontentsline{toc}{chapter}{LỜI NÓI ĐẦU}

\section*{1. Sự cần thiết của đề tài}
\addcontentsline{toc}{section}{1. Sự cần thiết của đề tài}

Hoạt động thể dục thể thao đang trở thành một nhu cầu thường xuyên của nhiều nhóm người dùng, đặc biệt là sinh viên, nhân viên văn phòng và những người muốn tự rèn luyện tại nhà. Tuy nhiên, việc bắt đầu và duy trì tập luyện không đơn giản nếu người dùng thiếu lộ trình, thiếu công cụ theo dõi hoặc không biết lựa chọn bài tập phù hợp với thể trạng hiện tại.


Trong thực tế, mỗi người dùng có mục tiêu, quỹ thời gian, thiết bị và tình trạng sức khỏe khác nhau. Một danh sách bài tập cố định chỉ hỗ trợ ở mức tham khảo, chưa giải quyết được yêu cầu cá nhân hóa. Vì vậy, đề tài xây dựng website tập luyện cá nhân hóa nhằm kết hợp dữ liệu hồ sơ sức khỏe, kho bài tập, kế hoạch luyện tập, thống kê tiến trình và chatbot hỗ trợ trong cùng một hệ thống.


Điểm trọng tâm của đề tài không phải là thay thế huấn luyện viên chuyên nghiệp, mà là tạo một môi trường web giúp người dùng phổ thông tiếp cận bài tập có chọn lọc hơn. Hệ thống cho phép đăng nhập, khai báo sức khỏe, xem bài tập, nhận gợi ý, tự lập kế hoạch, lưu tiến trình và nhận phản hồi từ trợ lý ảo. Các chức năng này được triển khai dựa trên mã nguồn thực tế của project Remix, Prisma, Auth0 và chatbot nhúng trên giao diện.


\section*{2. Mục tiêu và nhiệm vụ nghiên cứu}
\addcontentsline{toc}{section}{2. Mục tiêu và nhiệm vụ nghiên cứu}

Mục tiêu tổng quát là xây dựng một website hỗ trợ tập luyện cá nhân hóa theo thể trạng và nhu cầu của người dùng, đồng thời trình bày đầy đủ quá trình khảo sát, phân tích, thiết kế, triển khai và đánh giá sản phẩm trong phạm vi đồ án tốt nghiệp.


\begin{itemize}

    \item Khảo sát nhu cầu tập luyện tại nhà, nhận diện khó khăn của người dùng và xác định nhóm chức năng cần ưu tiên.

    \item Phân tích yêu cầu chức năng, yêu cầu phi chức năng, tác nhân sử dụng, luồng nghiệp vụ và tiêu chí chấp nhận của hệ thống.

    \item Thiết kế kiến trúc website, mô hình dữ liệu, luồng xử lý, nhóm API, Use Case, biểu đồ trình tự và nguyên tắc giao diện.

    \item Triển khai các chức năng chính từ mã nguồn: xác thực Auth0, quản lý hồ sơ sức khỏe, lọc bài tập, lập kế hoạch, ghi nhận tiến trình, thống kê và quản trị dữ liệu.

    \item Tích hợp chatbot hỗ trợ hỏi đáp về tập luyện, đồng thời nêu giới hạn an toàn để hệ thống không đóng vai trò chẩn đoán hoặc điều trị y tế.

    \item Kiểm thử các luồng chính và đánh giá ưu điểm, hạn chế, hướng phát triển của phiên bản thử nghiệm.

\end{itemize}

\section*{3. Đối tượng, phạm vi và phương pháp thực hiện}
\addcontentsline{toc}{section}{3. Đối tượng, phạm vi và phương pháp thực hiện}

Đối tượng nghiên cứu là quy trình người dùng tương tác với website tập luyện cá nhân hóa, bao gồm đăng nhập, cập nhật thông tin sức khỏe, xem bài tập, nhận gợi ý, lập kế hoạch, thực hiện bài tập, lưu tiến trình, xem thống kê và trao đổi với chatbot.


Phạm vi thực hiện tập trung vào phiên bản website thử nghiệm. Hệ thống không đưa ra khuyến nghị y khoa chuyên sâu, không kết nối thiết bị đeo thông minh trong phiên bản hiện tại và chưa triển khai cơ chế thanh toán hoặc gói huấn luyện trả phí. Các dữ liệu sức khỏe được dùng chủ yếu để lọc và sắp xếp bài tập ở mức phần mềm đồ án.


Phương pháp thực hiện gồm khảo sát nhu cầu, phân tích nghiệp vụ, thiết kế hệ thống theo hướng Use Case, xây dựng cơ sở dữ liệu, đọc mã nguồn triển khai, kiểm thử thủ công theo kịch bản và đối chiếu kết quả với yêu cầu đã đặt ra.


\section*{4. Bố cục báo cáo}
\addcontentsline{toc}{section}{4. Bố cục báo cáo}

\begin{itemize}

    \item Chương 1: CƠ SỞ LÝ THUYẾT

    \item Chương 2: KHẢO SÁT VÀ PHÂN TÍCH HỆ THỐNG

    \item Chương 3: THIẾT KẾ HỆ THỐNG

    \item Chương 4: TRIỂN KHAI

\end{itemize}
